\section{Introdução}\label{sec:introduction}

A preparação de um trabalho acadêmico envolve duas responsabilidades distintas: produzir conteúdo científico consistente e apresentar esse conteúdo de acordo com as regras aplicáveis. Em documentos extensos, ajustes manuais de margens, títulos, listas, citações e elementos pré-textuais aumentam o risco de inconsistência e dificultam a repetição do processo. Sistemas de preparação de documentos como o \LaTeX{} permitem separar conteúdo e apresentação e automatizar parte desse trabalho \parencite{lamport1994}.

Na Universidade Federal do Ceará (UFC), a normalização acadêmica combina normas técnicas vigentes e orientações institucionais publicadas pelo Sistema de Bibliotecas \parencite{ufc2026normalizacao}. Nesse contexto, uma classe documental pode reduzir tarefas repetitivas, mas não elimina a necessidade de revisar o texto, conferir as fontes consultadas e verificar os requisitos específicos do curso ou programa.

Este TCC é um exemplo didático. O tema escolhido permite mostrar, em um documento curto, como organizar um trabalho real e como usar os principais recursos do \texttt{abntexto-ufc}. Explicações detalhadas da API permanecem na documentação do projeto; o corpo deste trabalho segue a estrutura acadêmica usual.

\subsection{Problema}

Mesmo quando um modelo automatiza a apresentação, erros podem surgir por configuração inadequada, conteúdo incompleto ou diferenças entre o fonte e o PDF final. O problema considerado neste estudo é como organizar um fluxo simples capaz de compilar e verificar automaticamente propriedades básicas de um documento acadêmico antes da revisão humana.

\subsection{Objetivos}

O objetivo geral é descrever um fluxo demonstrativo de validação automatizada de documentos acadêmicos produzidos em \LaTeX{}.

Como objetivos específicos, pretende-se:

\begin{ufclettereditems}
  \item organizar o documento em etapas reprodutíveis de preparação, compilação e verificação;
  \item observar propriedades do PDF que podem ser verificadas automaticamente;
  \item registrar resultados de uma execução sintética do fluxo;
  \item destacar quais decisões continuam dependentes de revisão humana.
\end{ufclettereditems}

\subsection{Organização do trabalho}

Além desta introdução, a Seção~\ref{sec:background} apresenta os conceitos necessários; a Seção~\ref{sec:methodology} descreve o procedimento demonstrativo; a Seção~\ref{sec:results} reúne os resultados sintéticos; e a Seção~\ref{sec:conclusion} apresenta as considerações finais.
